\documentclass[11pt]{article}

\usepackage[margin=1in]{geometry}
\usepackage{amsmath,amssymb}
\usepackage{booktabs}
\usepackage[T1]{fontenc}
\usepackage[utf8]{inputenc}
\usepackage{lmodern}
\usepackage{microtype}
\usepackage{hyperref}
\usepackage{enumitem}
\usepackage{tabularx}

\hypersetup{
  colorlinks=true,
  linkcolor=blue,
  urlcolor=blue
}

\newcommand{\ThetaSet}{\Theta(\tau)}
\newcommand{\Varhat}{\widehat{\operatorname{Var}}}
\newcommand{\Covhat}{\widehat{\operatorname{Cov}}}

\title{The Current \texttt{scripts-paper} Workflow Explained}
\author{}
\date{18 July 2026}

\begin{document}

\maketitle

\begin{abstract}
This note describes the current source in \path{scripts-paper/}, including its
current fail-closed numerical policy and its present execution failure. The
sole production entrypoint, \path{scripts-paper/run_pipeline.R}, has 49
explicit \texttt{paper\_source\_once()} calls. Those calls enter a transitive
graph of 194 reachable production R modules. The tree currently contains 297
R files in total: 194 reachable production modules, 100 test modules, and
three documented inactive or test-support modules.

The workflow builds fitted stochastic-discount-factor (SDF) proxies, estimates
a consumption-growth mean equation, and is designed to map several
log-variance estimators over the mean equation's identified set. The current
backend, however, has no global solver for the positive-slack indefinite
quadratically constrained problems. It therefore marks every positive-slack
news-coefficient profile as \texttt{unreliable}. Only the closed-form
$\tau=0$ mean-equation point remains certified. The live tau-star search
therefore returns a typed \texttt{unresolved} status with lower floor zero and
no certified upper transition. The runner then stops in the structural-table
caption because that consumer still requests the removed scalar tau-star
field. A later selector-provenance mismatch also remains latent in PPML
coverage. Consequently, the current source does not produce a complete run,
final route record, or current set of downstream empirical results.
\end{abstract}

\tableofcontents

\section{Snapshot and scope}

This document describes the shared working tree as inspected on 18 July 2026.
That tree contains changes beyond \texttt{HEAD} commit
\texttt{cb7c8c6cad920db3fc6d718f9439fbb09f7f3465}, including the new modules
\path{support/data_transformations.R} and
\path{mean_equation/figures/region_availability.R}, and the new typed
tau helpers \path{support/identification/tau_formatting.R} and
\path{support/identification/interval_tables.R}. The
current source inventory is 297 R files: 194 reachable production modules,
100 test modules, and three explicitly inactive or test-support modules. The
companion
\path{docs/run_all_scripts_results.tex} reports the result of executing an
isolated snapshot of this source and distinguishes completed upstream
calculations from unavailable downstream results.

The package and paper layers remain separate:

\begin{itemize}[leftmargin=1.5em]
  \item \path{R/}, \path{tests/}, \path{man/}, \path{data/}, and \path{inst/}
    form the reusable R package.
  \item \path{scripts-paper/} is a self-contained paper workflow. It may call
    exported package functions, but paper-specific orchestration, solvers,
    diagnostics, tables, and figures live under \path{scripts-paper/}.
\end{itemize}

\section{Source graph and runtime architecture}

\subsection{One entrypoint and a load-once graph}

The runner first loads \path{config/paths.R} with one ordinary
\texttt{source()} call, because the path and load-once machinery does not yet
exist. Every later source edge uses \texttt{paper\_source\_once()}. The loader
normalizes paths, records loaded modules, rejects circular source chains, and
unregisters a module whose evaluation fails. A module can declare its own
dependencies, which is why the 49 explicit runner calls expand to 194
reachable production modules.

All modules evaluate in one shared global environment. That convention makes
execution order important, but the load-once registry prevents accidental
double evaluation. Topology tests reject direct
\texttt{source(paper\_path(...))} calls and unreachable production modules.
The current topology suite passes for all 297 R files and the 63-record
artifact manifest.

\subsection{Major layers}

\begin{table}[ht]
\centering
\small
\begin{tabularx}{\textwidth}{@{}lX@{}}
\toprule
Layer & Current responsibility \\
\midrule
\path{config/} &
  Analysis, reporting, diagnostics, artifact, and decision contracts. \\
\path{support/} &
  Runtime loading, typed artifacts, shared ACM input, data transformations,
  statistics, identification adapters, numerical search, and \LaTeX{} helpers. \\
\path{data_preparation/} &
  SDF proxies, consumption log growth, yield volatility, asset-return PCs,
  and expected-SDF/news PCs. \\
\path{mean_equation/} &
  OLS, the Lewbel system, point and positive-slack calculations, bootstrap,
  variance shares, tables, diagnostics, and region figures. \\
\path{log_variance/} &
  Estimator-neutral mapping plus log-OLS, PPML, Harvey, gated LAD, inference,
  diagnostics, panels, and figures. \\
\path{variance_bounds/} and \path{reports/} &
  Approximation-error summaries and descriptive reporting, both scheduled
  late in the runner. \\
\bottomrule
\end{tabularx}
\caption{The current paper-pipeline layers.}
\end{table}

\subsection{Configuration ownership}

Scientific and numerical settings are centralized rather than repeated:

\begin{sloppypar}
\begin{itemize}[leftmargin=1.5em]
  \item \path{config/analysis_contract.R} owns dimensions, tau grids,
    inference search controls, variance-share settings, and input contracts.
  \item \path{config/analysis.R} derives dates, maturity grids, PC dimensions,
    bootstrap budgets, worker counts, and the maintained
    \texttt{impose\_beta2r\_null} restriction.
  \item \path{config/reporting.R} and \path{config/diagnostics.R} own reporting
    and heteroskedasticity-diagnostic policies.
  \item \path{log_variance/estimators/controls.R} and
    \path{log_variance/diagnostics/protocols.R} own log-variance search and
    diagnostic protocols.
\end{itemize}
\end{sloppypar}

Three \texttt{HETID\_*} variables affect execution:
\path{HETID_BOOT_REPS}, \path{HETID_BOOT_CORES}, and
\path{HETID_ASSERT_EQUIV}. \path{FRED_API_KEY} separately changes the FRED
transport. The runner temporarily selects \texttt{L'Ecuyer-CMRG} around the
forked log-variance bootstrap and restores the caller's RNG kind. It does not
restore the previous \texttt{.Random.seed}.

\section{Artifact contract and startup lifecycle}

\begin{sloppypar}
\path{config/artifact_manifest_data.R} owns a 63-row artifact manifest:
53 \texttt{required}, six \texttt{conditional\_lad}, and four
\texttt{conditional\_egarch}. \path{config/artifacts.R} is the lookup facade;
\path{config/artifact_lifecycle.R} owns conditional cleanup and route-state
validation. CSV and RDS outputs use the typed serialization layer in
\path{support/artifacts/typed_artifacts.R}.
\end{sloppypar}

At startup the runner:

\begin{enumerate}[leftmargin=1.5em]
  \item loads artifact and analysis contracts;
  \item creates output directories;
  \item deletes all ten conditional LAD and EGARCH paths;
  \item loads and validates one quarterly ACM object for reuse by SDF
    construction and variance-bound calculations; and
  \item applies the FRED download patch before requesting consumption data.
\end{enumerate}

A successful run would finish by validating conditional branch presence and
writing \path{output/state/conditional_route_status.rds}. The current run
stops before that finalization step. Thus ``63 declared artifacts'' must not
be confused with a current, verified materialization of 53 or 59 artifacts.

\section{Data construction}

\subsection{ACM and SDF proxy panels}

The runner loads quarterly ACM yields and term premia once. For maturity $i$
and the three-month news step $s$, the package constructs the expected log
one-period price proxy $\hat n(i,t)$ and price revision
\[
  \Delta_{t+1}p_{t+i}^{(1)}
  = \hat n(i-s,t+1)-\hat n(i,t).
\]
The paper then calls exported package functions to build:

\begin{itemize}[leftmargin=1.5em]
  \item an in-sample fitted expected-SDF proxy; and
  \item a centered second-order SDF-news proxy
    \[
    \exp\{\hat n(i,t)\}
      \left[\Delta p+\tfrac12(\Delta p)^2\right]-B_i.
    \]
\end{itemize}

These labels do not make the series exact observed conditional expectations
or revisions. Their principal components are generated regressors.

\subsection{Consumption growth}

\path{data_preparation/build_consumption_growth.R} obtains quarterly real
personal consumption expenditures, FRED series \texttt{PCECC96}, and calls
\path{support/data_transformations.R}. The current transform is
\[
  \Delta c_t
  =100\{\log(\text{PCECC96}_t)-\log(\text{PCECC96}_{t-1})\},
\]
quarterly log growth in percent, not annualized. The helper rejects
nonfinite or nonpositive observed levels. This replaces the arithmetic
$100(P_t/P_{t-1}-1)$ transform used by the older tracked outputs.

\subsection{The three PC families}

The workflow uses three distinct families:

\begin{itemize}[leftmargin=1.5em]
  \item $PC_E$: three PCs from the lagged fitted expected-SDF panel;
  \item $PC_N$: three PCs from the second-order SDF-news proxy panel; and
  \item $PC_R$: four precomputed PCs of a cross-section of financial asset
    returns, relabeled as \texttt{l.pc1} through \texttt{l.pc4}.
\end{itemize}

The asset-return PCs are never PCs of yields. Yields and term premia enter the
SDF proxy construction, not the $PC_R$ construction. An unlagged expected-SDF
PCA supplies only a sign-alignment reference; a separate lagged fit supplies
the structural $PC_E$ scores.

\section{Mean equation}

\subsection{Model and maintained restriction}

The mean equation is
\[
  \Delta c_{t+1}
  = b_0+PC_{E,t}'b_E+PC_{N,t+1}'\theta+\varepsilon_{t+1}.
\]
Regressing $\Delta c_{t+1}$ on $(1,PC_{E,t})$ gives $W_1$. The current
specification maintains $\beta_2^R=0$ for the reduced-form projection of
$PC_N$ on that specified common design, so $W_2=PC_N$. This is a
proxy-level identifying restriction. It is not implied by unconditional PCA
orthogonality or by the fitted SDF formulas.

The single instrument is de-meaned five-year realized yield volatility,
$Z=\texttt{y60\_vol}-\overline{\texttt{y60\_vol}}$. With one instrument,
the weight matrix is fixed and there is no weight optimization.

\subsection{Lewbel quadratic system}

The package computes the identification moments and returns one quadratic
constraint for each news component:
\[
  \ThetaSet
  =\{\theta:\theta'A_i(\tau)\theta+b_i(\tau)'\theta+c_i(\tau)\leq0,
    \ i=1,2,3\}.
\]
At $\tau=0$ the constraints collapse to a linear system $Q\theta=L$. When
$Q$ has full column rank and is consistent,
\texttt{solve\_point\_identification()} returns the closed-form point. This
is the only currently certified news-coefficient result.

\subsection{Current positive-slack certification policy}

For $\tau>0$, the feasible set is defined by indefinite quadratic
inequalities. The available SLSQP backend supplies local feasible candidates,
not global extrema. Current code makes that limitation operational:

\begin{itemize}[leftmargin=1.5em]
  \item \path{support/identification/profile_solver_core.R} assigns
    \texttt{valid=FALSE} to every finite local coordinate or linear-functional
    candidate.
  \item \path{support/identification/bound_search_classifier.R} also assigns
    \texttt{valid=FALSE} to an apparent box-growth infinity.
  \item \texttt{coef\_interval\_tables()} consequently marks every
    positive-slack $\theta$ row \texttt{unreliable}.
  \item Under the maintained $\beta_2^R=0$ restriction, $b_0$ and $b_E$ are
    separately known constants. Their positive-slack table cells remain
    blank rather than displaying repeated values.
\end{itemize}

Warm continuation from the $\tau=0$ point remains a numerical contradiction
diagnostic. It carries feasible starts forward, but it cannot upgrade a row:
replacement is guarded by \texttt{status == "bounded"}, which the current
local backend cannot produce for a positive-slack news coefficient.

\subsection{The typed \texorpdfstring{$\tau^\ast$}{tau-star} status}

\texttt{eval\_width\_at\_tau()} classifies an invalid solve as
\texttt{unreliable} before considering boundedness. Because all
positive-slack sides are invalid, the current tau-star sweep has one certified
bounded point, $\tau=0$, but neither a certified positive-slack bounded point
nor a certified unbounded point. \texttt{tau\_star\_fixed()} now returns a
typed result:
\[
  \texttt{status=unresolved},\qquad
  \texttt{lower=0},\qquad
  \texttt{upper=NA}.
\]
This is a certified floor at zero and no certified transition; it is not a
capped estimate. The former scalar \texttt{tau\_star} and Boolean
\texttt{tau\_star\_capped} fields no longer exist.

The 200-draw verification produced \texttt{unresolved} in every draw, a
lower-bound band $[0,0]$, and zero draws whose certified lower floor reached
the $0.05$ baseline. These are status diagnostics, not evidence that a
positive-slack set is bounded or unbounded.

\subsection{Bootstrap and rendering consequences}

The moving-block bootstrap still re-estimates the closed-form $\tau=0$ point,
so robust point intervals remain available. For positive slack, draw endpoints
are retained only when their status is \texttt{bounded}. There are therefore
no finite news-endpoint draws, and the symmetric joint-normal endpoint
calibration is suppressed. The implemented common-critical-value calibration
is bespoke; it is not Stoye's CI2 or CI3. If endpoint correlation were
unavailable, it would fall back to the Imbens--Manski interpolation.

The structural renderer is designed to print \texttt{unreliable} for each
positive-slack news coefficient, and variance-share rows derived from those
profiles are likewise unavailable. The current runner does not reach either
completed table, however, because its structural caption still requests the
removed scalar tau-star field; Section~\ref{sec:first-failure} gives the exact
failure. If execution reached the new region gate, the two mean-region SVGs
would contain explicit ``Unavailable'' placeholders instead of uncertified
geometry.

No local grid-and-polish result is now exposed as a block-share range. A
constant expected-SDF block can be certified under the maintained
$\beta_2^R=0$ restriction when the $\tau=0$ point certifies nonemptiness;
every nonconstant block range is withheld. The component-share mapper also
fails closed on disconnected projections. A coordinate interval that strictly
straddles zero does not prove that zero is feasible for an indefinite
quadratic system. Unless an endpoint is literally zero, such an otherwise
bounded straddle is downgraded to \texttt{unreliable} rather than assigned a
zero minimum share.

\section{Log-variance layer}

\subsection{Intended estimands}

For a candidate $\theta$, define
\[
  \varepsilon_t(\theta)=W_{1,t}-W_{2,t}'\theta.
\]
The intended log-variance maps use four lagged asset-return PCs:

\begin{align*}
 E[\log\varepsilon_t(\theta)^2\mid PC_{R,t-1}]
   &=\vartheta_0^L(\theta)+PC_{R,t-1}'\vartheta_R^L(\theta),\\
 E[\varepsilon_t(\theta)^2\mid PC_{R,t-1}]
   &=\exp\{\vartheta_0^V(\theta)+PC_{R,t-1}'\vartheta_R^V(\theta)\}.
\end{align*}

Log-OLS and LAD use the log-squared target; PPML and Harvey use the
level-scale second moment. The estimator-neutral engine is designed to scan a
certified mean-equation box, evaluate feasible points, and locally polish
coordinate extrema. Such attained values would be inner approximations, not
global projection hulls.

\subsection{Fail-closed propagation}

\path{log_variance/engine/api.R} checks the upstream mean box before scanning.
When any required mean side is \texttt{unreliable}, it returns all-NA
log-variance endpoints with \texttt{unreliable} status. It does not run the
Morton grid selector. This is the correct fail-closed response to the current
mean-set policy.

This behavior is verified at the module level and was observed in an earlier
isolated snapshot. The final live runner now stops at the structural-table
caption before log-OLS is loaded, so it produces no current log-variance
diagnostic, residual-crossing count, or positive-slack range.

\subsection{Current first hard stop}
\label{sec:first-failure}

After data preparation, mean estimation, and the mean bootstrap, runner
line~38 loads the structural-table renderer. The estimated object stores the
new fields
\texttt{tau\_star\_status}, \texttt{tau\_star\_lower}, and
\texttt{tau\_star\_upper}. Its caption still passes the removed
\texttt{set\_id\_mean\_eq\$tau\_star} field to
\texttt{tau\_transition\_range\_label()}. That value is \texttt{NULL}, so the
formatter's \texttt{switch()} stops with:
\begin{quote}
\texttt{EXPR must be a length 1 vector}
\end{quote}

The call path is
\begin{quote}
\path{run_pipeline.R:38}\\
$\longrightarrow$
\path{mean_equation/tables/render_structural_equation_table.R}\\
$\longrightarrow$
\path{mean_equation/tables/structural_equation_caption.R}\\
$\longrightarrow$ \texttt{tau\_transition\_range\_label()}.
\end{quote}

There is also a later, latent PPML mismatch. In the immediately preceding
snapshot, which reached PPML, the fail-closed engine returned completed
records without selector diagnostics and coverage stopped with
\texttt{coverage selector provenance is absent from a completed engine run}.
That remains a source-level inconsistency, but it is not the first failure of
the final live tree.

\subsection{Modules not reached}

The hard stop occurs before:

\begin{itemize}[leftmargin=1.5em]
  \item completed structural-equation tables and all variance-share outputs;
  \item log-OLS and the mean bounds-by-tau figure;
  \item PPML estimation, standard errors, and tables;
  \item Harvey estimation;
  \item joint-null and joint-GMM diagnostics;
  \item the residual-dynamics gate and EGARCH router;
  \item the gated LAD estimator;
  \item log-variance set bootstrap, panels, and fitted-volatility figures;
  \item mean-region placeholder rendering;
  \item heteroskedasticity table rendering, variance-bound rendering,
    descriptive reports, and final conditional-route validation.
\end{itemize}

The joint-null module contains a truthful unavailable-state branch for
uncertified boxes, and the fitted-volatility and region renderers contain
unavailable placeholders. They describe intended fail-closed behavior if
reached; they are not current-run artifacts.

\section{Conditional branches}

LAD and EGARCH use different controls:

\begin{sloppypar}
\begin{itemize}[leftmargin=1.5em]
  \item LAD reads the tracked DCF record
    \path{config/decisions/lad.dcf}, currently approved for
    \texttt{quantreg} 6.1.
  \item EGARCH validates a hash-bound R decision against the regenerated
    dynamics gate, fixed plan digest, upstream-plan digest, and prompt hashes.
    No EGARCH estimator is currently wired.
\end{itemize}
\end{sloppypar}

These controls remain part of the architecture, but the current runner stops
before either route executes. Startup still removes their ten conditional
artifacts. Because final route validation is not reached, there is no current
authoritative claim that LAD was requested or ran, nor that a 59-file route
was materialized.

\section{Current inconsistencies and how to read the tree}

\begin{table}[ht]
\centering
\small
\begin{tabularx}{\textwidth}{@{}p{0.28\textwidth}X@{}}
\toprule
Current symptom & Correct interpretation \\
\midrule
Tracked tables contain positive-slack ranges and old consumption statistics &
They were generated by an older arithmetic-growth, locally certified-box
workflow and are stale relative to current source. \\
Tau-star estimation is typed, but the structural caption requests the removed
scalar field &
The live status is \texttt{unresolved} with lower floor zero and no upper
transition. Passing the absent field to the formatter is the current first
hard failure. \\
Variance-share output is scheduled after the current failure &
If reached, the mapper would certify only a constant block on a certified
nonempty set, withhold every nonconstant block range, and downgrade a strict
component straddle without a literal zero endpoint. \\
Manifest declares 63 paths &
This is a declaration, not evidence of a completed current run. \\
Downstream modules contain unavailable-state logic &
The structural-caption error occurs earlier, so the full runner does not
reach those handlers. A second selector-provenance mismatch remains latent in
PPML coverage. \\
\bottomrule
\end{tabularx}
\caption{Important distinctions in the present working tree.}
\end{table}

\section{Reproducing the present behavior}

From the package root, the full command remains:
\begin{verbatim}
Rscript scripts-paper/run_pipeline.R
\end{verbatim}

For the verification described here, the isolated run used:
\begin{verbatim}
HETID_BOOT_REPS=200 HETID_BOOT_CORES=2 HETID_ASSERT_EQUIV=1 \
  Rscript scripts-paper/run_pipeline.R
\end{verbatim}

It completed data preparation, mean estimation, and all 200 mean-bootstrap
draws, then stopped while building the structural-table caption with the exact
message quoted above. A successful exit should not be expected until the
caption consumes the typed tau-star fields. Reaching the later log-variance
path would still expose the selector-free PPML coverage mismatch described
above.

\section{Bottom line}

The current design is deliberately conservative about global numerical
claims. It correctly refuses to promote local positive-slack QCQP candidates
to identified-set bounds. That policy leaves the closed-form $\tau=0$ mean
point as the only certified structural result. The current runner first fails
because its structural caption has not been migrated to the typed tau-star
object. A later PPML coverage path also treats a selector-free fail-closed
result as an error. Until both control-flow mismatches are repaired and a
complete run regenerates the manifest, older positive-slack, log-variance,
dynamics, and route artifacts are historical rather than current results.

\section*{References}

\begin{itemize}[leftmargin=1.5em]
  \item Imbens, G. W., and Manski, C. F. (2004). Confidence intervals for
    partially identified parameters. \emph{Econometrica}, 72(6), 1845--1857.
  \item Lewbel, A. (2012). Using heteroscedasticity to identify and estimate
    mismeasured and endogenous regressor models.
    \emph{Journal of Business \& Economic Statistics}, 30(1), 67--80.
  \item Stoye, J. (2009). More on confidence intervals for partially
    identified parameters. \emph{Econometrica}, 77(4), 1299--1315.
\end{itemize}

\end{document}
